\section{The Thurstonian Model of a Field}
\label{sec:model}

Consider a universe of $n$ assets. Each asset has a latent ability $a_i\in\R$ and
a noisy \emph{performance}
\begin{equation}
  X_i = a_i + \varepsilon_i, \qquad \varepsilon = (\varepsilon_1,\dots,\varepsilon_n)
  \sim S,
  \label{eq:performance}
\end{equation}
where $S$ is any centered joint law---the \emph{simulation} that drives the
race---and we adopt the convention that the \emph{minimum} performance wins, so
that a smaller ability denotes a stronger competitor. The Gaussian
$S = \mathcal{N}(0, \Sigma_\varepsilon)$ is the workhorse, and the only case in
which the calibration and smoothness results below are closed-form; but the
construction---and the implementation---treat $S$ as a black-box sampler, so a
copula with genuine tail dependence, a fat-tailed or skewed law, or a structured
generator are equally admissible (\Cref{prop:tail,rem:anysim}). The \emph{state price} (winning
probability) of asset $i$ is
\begin{equation}
  p_i \;=\; \Prob\!\left( X_i = \min_{k} X_k \right).
  \label{eq:stateprice}
\end{equation}
When the performances are independent ($\Sigma_\varepsilon$ diagonal) the
$p_i$ are determined by the abilities alone through the distribution of the first
order statistic, and $p$ is a smooth, monotone function of $a$
\parencite{cotton2021}. Two maps are then of interest:
\begin{align}
  \text{forward:}\quad & a \;\longmapsto\; p, \label{eq:forward}\\
  \text{inverse:}\quad & p \;\longmapsto\; a. \label{eq:inverse}
\end{align}
The inverse \eqref{eq:inverse}---the \emph{Fast Ability Transform}---recovers
abilities (up to an additive constant) from a target set of winning
probabilities, and is the engine of the construction below.

\section{Thurstone Portfolios via the Ability Tilt}
\label{sec:method}

We treat portfolio \emph{weights} as winning probabilities and \emph{assets} as
competitors. The construction uses two simulations---a \emph{reference} law and a
\emph{target} law. We first \emph{calibrate} latent abilities so that the race
under the reference law reproduces a target allocation $w^{\mathrm{target}}$ (a
benchmark such as capitalization weights); we then re-run the race under the
target law and read off the winning probabilities as the tilted weights. In the
Gaussian workhorse the two laws are two correlation matrices---a reference
$C_{\mathrm{calib}}$ and the estimated return correlation $C_{\mathrm{tilt}}$---and
the tilt is the portfolio's response to $C_{\mathrm{tilt}} - C_{\mathrm{calib}}$
(\Cref{prop:fixedpoint}); but any pair of centered samplers may be used, with the
target law in particular carrying dependence (tail dependence included) that the
reference does not.

\begin{algorithm}[h]
\caption{Ability tilt (Thurstone portfolio)}
\label{alg:tilt}
\begin{algorithmic}[1]
\Require target weights $w^{\mathrm{target}}\in\Delta$; reference sampler
         $S_{\mathrm{calib}}$; target sampler $S_{\mathrm{tilt}}$ (Gaussian
         $\mathcal{N}(\,\cdot\,, C_{\mathrm{tilt}})$ by default).
\State $a \gets \textsc{Calibrate}(w^{\mathrm{target}}, S_{\mathrm{calib}})$
       \Comment{abilities s.t.\ the race under $S_{\mathrm{calib}}$ yields $w^{\mathrm{target}}$}
\State Draw $X^{(1)},\dots,X^{(M)} \sim S_{\mathrm{tilt}}(a)$ from \emph{fixed} seeds
       \Comment{Gaussian: $X^{(m)} = a + C_{\mathrm{tilt}}^{1/2} Z^{(m)}$}
\State $w_i \gets \dfrac{1}{M}\sum_{m=1}^{M}
       \mathbf{1}\!\left[\, i = \argminop_k X^{(m)}_k \,\right]$
       \Comment{win frequency}
\State \Return $w$
\end{algorithmic}
\end{algorithm}

Two choices of the reference correlation give the practical flavours:
\begin{itemize}
  \item \textbf{Diagonal} ($C_{\mathrm{calib}} = I$). Calibration is the exact
        order-statistic inverse of \Cref{sec:model} (the Fast Ability
        Transform), and the tilt then injects the estimated correlation. The
        target may be any benchmark---the variance-only diagonal portfolio,
        capitalization weights, or equal weight.
  \item \textbf{One-factor market}
        ($C_{\mathrm{calib}} = \beta\beta^\top + \mathrm{diag}(1-\beta^2)$). The
        reference is a single market factor with loadings $\beta$; calibration
        reproduces the benchmark \emph{given} that market correlation, so the
        tilt's deviation is driven by the residual (non-market) correlation in
        $C_{\mathrm{tilt}}$.
\end{itemize}

\noindent In the Gaussian workhorse a scalar $\phi \in [0,1]$ sets
$C_{\mathrm{tilt}} = (1-\phi)\,C_{\mathrm{calib}} + \phi\,\widehat{C}$, dialling
the estimated correlation $\widehat{C}$ in continuously: $\phi = 0$ reproduces the
benchmark (\Cref{prop:fixedpoint}) and $\phi = 1$ uses the full estimate. The
construction never inverts a covariance---it needs only a valid correlation for
sampling, repaired if necessary to the nearest correlation matrix, so it tolerates
ill-conditioned or singular estimates. Its feasibility, fixed-point, redundancy,
monotonicity, and objective properties are established in \Cref{sec:properties}.

\section{Computation}
\label{sec:computation}

\paragraph{Calibration (deterministic).}
The calibration step---the inverse map from target weights to abilities---is
computed deterministically with the lattice machinery of the \texttt{thurstone}
package, which builds the distribution of the first order statistic by convolving
competitor densities and tracking conditional multiplicity for ties
\parencite{cotton2021}. For a diagonal reference this is the exact
order-statistic inverse \eqref{eq:inverse}, with no Monte-Carlo error and cost
linear in $n$. For the one-factor reference the assets are conditionally
independent given the market factor, so the same lattice race is evaluated at a
few Gauss--Hermite nodes and integrated over the factor; the inverse is then a
damped fixed-point on this forward map. This is the only place quadrature enters.

\paragraph{Tilt (Monte-Carlo).}
The race under a general target law has no tractable analytic form, so it is
evaluated by sampling. The \texttt{allocation} package factors the race into a
\emph{draw}---the sampler---and a \emph{scoring} step, the argmin win-frequency
core, so the sampler is a plug-in: the Gaussian default recolours a fixed
standard-normal ensemble by $C_{\mathrm{tilt}}^{1/2}$ (a quasi-Monte-Carlo
ensemble may be substituted to sharpen the win-frequency \emph{level}, though not
the step-to-step turnover of \Cref{sec:smoothness}); a Student-$t$ and a low-rank
factor race (for scale) ship alongside; and any user callable obeying the
fixed-seed contract supplies an arbitrary law. Tail consistency (\Cref{prop:tail}) follows by
driving the race with a downside-dependent sampler---or, within the Gaussian
workhorse, a downside (lower-partial-moment) covariance. This supersedes the
earlier \texttt{winningport} implementation of the tilt.

\paragraph{Online updates and scale.}
At rebalancing the correlation estimate moves slowly, and \Cref{sec:smoothness}
shows the weights move smoothly with it. We exploit this by holding the
standard-normal seeds fixed and \emph{transporting} the path ensemble to the new
correlation (recolouring by the updated square root), so turnover tracks genuine
correlation change rather than sampling noise. For large universes the dense
correlation is never formed: a low-rank factor model $C = BB^\top + D$ keeps both
the transport and the argmin streaming and linear in $n$, scaling to thousands of
assets. The covariance estimate itself is supplied by any external online
estimator, so the method interoperates with standard covariance-forecasting
tooling.

